\documentclass[a4paper,12pt]{article}
\usepackage{graphicx} % For inserting images
\usepackage[serbian]{babel} % Serbian language support
\usepackage[a4paper, margin=2.5cm]{geometry} % Improved page layout
\usepackage{fancyhdr} % Custom headers/footers
\usepackage{titlesec} % Better section titles
\usepackage{tocloft} % Better TOC spacing
\usepackage{biblatex} % Bibliography management
\usepackage{csquotes} % Recommended for biblatex
\usepackage[colorlinks=true, urlcolor=cyan, linkcolor=blue]{hyperref}
\usepackage[T1]{fontenc} % For đ


\setlength{\headheight}{14.5pt} % Fix fancyhdr warning

\pagestyle{fancy}
\fancyhf{} % Clear default header/footer
\fancyhead[L]{\small \leftmark} % Header left
\fancyfoot[C]{\thepage} % Page number in center footer

\titleformat{\section}{\large\bfseries}{\thesection}{1em}{}

\addbibresource{references.bib} % Reference file

\title{Računarstvo i društvo - pregled gradiva}
\author{Katarina Grbović}
\date{Novembar 2025}

\renewcommand{\contentsname}{Sadržaj}

\begin{document}

\maketitle
\newpage
\tableofcontents
\newpage
\section{Istorijski pregled razvoja računarstva}
\subsection{Pregled gradiva}
Razvoj računarstva možemo podeliti u tri celine: 
%\begin{itemize}
%    \item Razvoj fizičkih računara
%    \item Razvoj programskih jezika i softvera
%    \item Razvoj interneta
%\end{itemize}
Teško je definisati tačan početak razvoja računara, ali možemo posmatrati razvoj računara od nastanka bušenih kartica, koje predstavljaju prvi medijum za čuvanje podataka. Nakon toga beležimo nastanak elektronskih računara 1930-ih godina, gde su se instrukcije programirale umetanjem žica u odgovarajuće otvore. Već 1950-ih godina se razvija način davanja instrukcija računarima i tako se javljaju prvi programski jezici. Paralelno sa razvojem programskih jezika razvija se i novi hardver, koji je neuporedivo manji od prvih računara i neuporedivo moćniji. Već 1970-ih godina javljaju se prvi personalni računari, neuporedivi sa računarima od pre pola veka. Simultano, već od 19. veka razvija se i komunikacija prvim električnim telegrafom nastao 1830-ih godina. U međuvremenu nastaju i razvijaju se telefon, radio, televizija. 
1967. godine nastaje ARPANET, preteča današnjeg interneta.
\subsection{Predložene teme}
\begin{itemize}
    \item Request for Comments i kakav uticaj je imao RFC na razvoj interneta
    \begin{itemize}
        \item \url{https://www.rfc-editor.org/history/}
    \end{itemize}
\end{itemize}
\section{Etika}
\subsection{Pregled gradiva}
\textbf{Etika} je grana filozofije koja se bavi izučavanjem morala, racionalnom sistematizacijom ljudskih verovanja i ponašanja, kao i odbranom i preporučivanjem koncepata ispravnog i pogrešnog ponašanja.
Podoblasti etike su:
\begin{itemize}
    \item Metaetika
    \item Normativna etika
    \item Primenjena etika
\end{itemize}
Etičke teorije - dobra teorija etike omogućava pravljenje ubedljivih logičkih argumenata u vezi sa posmatranim problemom prema široj publici.
Podela etičkih teorija:
\begin{itemize}
    \item Subjektivni relativizam
    \item Kulturološki relativizam
    \item Teorija duhovnog komandovanja
    \item Etički egoizam
    \item Kantijanizanm
    \item Utilitarizam postupaka
    \item Utilitarizam pravila 
    \item Teorija socijalnog ugovora
    \item Etika vrline
\end{itemize}
Odnos etike i tehnološkog razvoja. Etika tehnologije.
\begin{itemize}
    \item Tehnicizam
    \item Optimizam
    \item Pesimizam
\end{itemize}
Računarska etika i pristupi računarskoj etici.
\subsection{Predložene teme}
\begin{itemize}
    \item Siva zona sajberbezbednosti
    \begin{itemize}
        \item \url{https://www.eccu.edu/blog/cybersecurity/the-ethics-of-cybersecurity-debating-the-gray-areas/}
    \end{itemize}
\end{itemize}

\section{Internet i spam poruke}
\subsection{Pregled gradiva}
Prvi spam telegraf je poslat 1864. godine, a prva spam poruka putem imejla 1978. godine. Do 2009. godine, spam poruke čine oko 8 posto svih poruka. Od 2009. količina spama kreće eksponencijalno da raste. Spam poruke imaju obično svrhu da na daleko jeftiniji način reklamiraju određeni proizvod u odnosu na bilo koji drugi medijum, a njihov sadržaj su upravo najčešće reklame, zatim eksplicitan sadržaj, finansijski sadržaj i prevara. Spameri najčešće pronalaze imejlove putem whois baza, veb stranica, ili pak nasumičnog generisanja imejl adresa. Najkorišćenija anti spam mera su upravo spam filteri.
Razlika između interneta i veba je da internet čini mrežnu infrastrukturu i internet protokole dok je WWW usluga definisana i koristi HTTP protokol i zavisi od interenta. Evolucija veba kreće već od 70-ih godina prošlog veka i veb se komercijalizuje devedesetih godina prošlog veka i kreće polako sa centralizacijom veba koja traje i dan danas. Sajtovi su do pre nekoliko godina bili uglavnom statički, odnosno server je servirao sadržaj od kog nije dalje zavisio. Međutim, sa rastom popularnosti jezika poput javaskripta, na mnogo većoj popularnosti dobijaju dinamički sajtovi.
Internet se koristi u najrazličitije svrhe. Elektronska pošta se smatra kao osnov komunikacije putem interneta, dok su socijalne mreže savremeni potomak. Pored ovoga, popularni su i onlajn kupovina, elektronsko bankarstvo i sl.
\subsection{Predložene teme}
\begin{itemize}
    \item Alati za presretanje paketa u internet saobraćaju i njihova uloga
    \begin{itemize}
        \item \url{https://www.ijfmr.com/papers/2023/6/8876.pdf}
    \end{itemize}
\end{itemize}
\section{Cenzura i sloboda govora, deca i neprikladan sadržaj}
\subsection{Pregled gradiva}
Cenzura je autoritarno ograničavanje ili zabrana javnog izražavanja mišljenja ili ideja, zbog opasnosti po vlast, društveni ili moralni poredak.
Pregled cenzure kroz istoriju, kako štamparska mašina utiče na cenzuru i kako se cenzura menja. Licenciranje i registracija stanica i kanala.
Vrste cenzure: 
\begin{itemize}
    \item Direktna
    \item Autocenzura
\end{itemize}
Na osnovu istraživanja, korisnici najviše vrše autocenzuru u oblastima svakodnevnog života i ličnog mišljenja. Mnogi autoritarni režimi imaju monopol nad internetom gde se vrši cenzura interneta. Iako je nemogućnost potpune cenzure interneta u svojoj osnovi dobra stvar, manjak cenzure utiče loše na decu koja koriste internet sve mlađa i mlađa. Deca su izložena, čak i uz najobičnniju pretragu, eksplicitnom sadržaju, samopovređivanju, pornografiji. 
Postoje mnogi tipovi neprikladnog sadržaja na koje dete može naići i samim tim postavlja se pitanje kako zaštititi decu od eksplicitnog i nedoličnog sadržaja. Primer zaštite deteta od necenzurisanog intereneta su internet filteri ili roditeljski nadzor, kao i otvorena i iskrena komunikacija sa decom i zajedničko istraživanje interneta.
\subsection{Predložene teme}
\begin{itemize}
    \item Nadzor i cenzura i kako tehnologija utiče na ljudska prava
    \begin{itemize}
        \item \url{https://www.coe.int/en/web/commissioner/-/highly-intrusive-spyware-threatens-the-essence-of-human-rights}
    \end{itemize}
\end{itemize}
\section{Internet prevare i zavisnost od interneta}
\subsection{Pregled gradiva}
Koliko prednosti Internet ima, toliko ima i mana. Javlja se velika zavisnost od interneta što postaje globalni problem koji sa sobom vuče brojne zdravstvene probleme. Između ostalog, jedna od brojnih loših strana interneta su upravo internet prevare. 
\begin{itemize}
    \item Krađa identiteta - korišćenje tuđih podataka zarad pristupa podacima te osobe
    \item Phishing - slanje poruka u kojima se urgentno traži da korisnik na neki način unese svoje podatke iza čega stoji prevara
    \item Nigerijska prevara - traži se od žrtve pomoć za transfer velike svote novca uz obećanje da žrtva dobija procenat
    \item Lažne ocene proizvoda i usluga - barem jedna trećina ocena na internetu su lažne
    \item Dezinformacije - Uzevši u obzir neverovatan broj internet stranica sa različitim sadržajem neophodno je primetiti i da je plasiranje i širenje dezinformacija na internetu veliki problem
    
\end{itemize}
Zavisnost od interneta je relativno nov fenomen i zahteva sopstvenu klasifikaciju i ima drugačije simptome od nekih opšte poznatih tipova zavisnosti. Neki od problema koji su direktna posledica zavisnosti od interneta su depresija, anksioznost, promene raspoloženja, insomnija, problemi sa vidom i slično. Simptomi zavisnosti od društvenih mreža su slični kao kod bolesti zavisnosti prouzrokovanih hemijskim supstancama. Zavisnost od internet igara takođe predstavlja jednu od podkategorija zavisnosti od interneta.
\subsection{Predložene teme}
\begin{itemize}
    \item Kako se boriti sa širenjem dezinformacija na internetu
    \begin{itemize}
        \item \url{https://www.brookings.edu/articles/how-to-combat-fake-news-and-disinformation/}
    \end{itemize}
\end{itemize}
\section{Bezbednost interneta}
\subsection{Pregled gradiva}

Hakovanje i istorija hakovanja datiraju još pre personalnih računara. Iako je zlatno doba hakovanja davno prošlo, savremeni internet korisnici se još uvek bore sa hakerima i hakerskim napadima. Neophodno je zaštititi svoje naloge od napadača postavljanjem kompleksnih šifara, ali opet treba biti obazriv od napada u okviru besplatne WiFi mreže jer haker može pristupiti našim kolačićima koje imamo u toj sesiji.
Virusi predstavljaju maliciozne procese koji se izvršavaju na računaru, a antivirusi su softveri zaduženi da ih detektuju.
Podela zlonamernog softvera:
\begin{itemize}
    \item Kompjuterski crvi
    \item Cross Site Scripting
    \item Rootkit
    \item Exploit
    \item DDoS
    \item Backdoor
    \item Spajver
    \item Botnets
\end{itemize}
Šifrovanje, rainbow tables, salt, enkripcija, heširanje su sve deo čuvanja i vektora napada kada je u pitanju krađa tuđih šifara.
\subsection{Predložene teme}
\begin{itemize}
    \item Istorija kompjuterskih virusa i njihova promena koja prati razvoj interneta u 21. veku
    \begin{itemize}
        \item \url{https://www.kaspersky.com/resource-center/threats/a-brief-history-of-computer-viruses-and-what-the-future-holds}
    \end{itemize}
\end{itemize}
\section{Privatnost na internetu}
\subsection{Pregled gradiva}
Šta je privatnost i kako je definišemo? Ključni pojam u raspravama o tome šta je privatnost je koncept pristupa. Da li mi svaki put svesno i voljno ostavljamo naše informacije na internetu, i da li smo uopšte svesni kada se to dešava protiv naše volje? Korisnici sve više brišu informacije o sebi i veća je upotreba privatnih naloga nego pre. 
Kolačići su male tekstualne datoteke koje čuvaju informacije o veb sajtovima. Delimo ih na:
\begin{itemize}
    \item Sesijske
    \item Trajne - oni se dalje dele na kolačiće prve i treće strane
\end{itemize}
Do skoro je korišćenje kolačića bilo implicitno, odnosno podrazumevalo se da korisnik treba sam da isključi opciju kolačića ukoliko to želi. Ovaj problem je sada zakonski regulisan te svi sajtovi moraju da obaveste korisnika da prikupljaju informacije o njima putem kolačića.
Koje to informacije o nama kolačići sve čuvaju?
\begin{itemize}
    \item Istorija pretraživanja, kupovine
    \item Istorija lokacije
    \item Informacije o uređajima
    \item Kada i gde je viđena prethodna reklama
    \item Linkovi na koje je korisnik kliknuo
\end{itemize}
Kako se možemo zaštititi od kolačića? 
\begin{itemize}
    \item Korišćenje pregledača u privatnom režimu rada
    \item Onemogućiti čuvanje kolačića u podešavanjima pregledača
    \item Brisanjem kolačića
\end{itemize}
Koliku kontrolu ima država nad svakodnevnim stanovništvom putem digitalnog nadzora, kako možemo uticati na zaštitu privatnosti? 
Neki od primera zaštitne privatnosti su enkripcija i VPN.
\subsection{Predložene teme}
\begin{itemize}
    \item Da li nas VPN stvarno štiti i koje informacije on prikuplja o nama?
    \begin{itemize}
        \item \url{https://allaboutcookies.org/are-vpns-safe}
    \end{itemize}
\end{itemize}
\end{document}
